\documentclass{article}
\usepackage{amsfonts,amsmath,amssymb,amsthm}
\usepackage{hyperref} % do NOT set [ocgcolorlinks] here!
\usepackage{nameref}
\usepackage{xcolor}

\newtheorem{lemma}{Lemma}

\newenvironment{lemmaproof}
  {\renewcommand{\qedsymbol}{$\blacksquare$}\begin{proof}}
  {\end{proof}}

\title{$\aleph_{0}$ Weekly Problem}
\author{Ravi Dayabhai}
\date{2025-11-30}

\begin{document}

\maketitle

\section*{Problem}

Let $a_n$ be the number of terms in the sequence $2^{1}, 2^{2}, \ldots, 2^{n}$ which begin with digit 1. Prove that
  \[
    \log{2} - \frac{1}{n} < \frac{a_{n}}{n} < \log{2},
  \]
and, hence,
  \[
    \lim_{n \to \infty} \frac{a_{n}}{n} = \log{2} \approx 0.30103. 
  \]

\section*{Solution}

  \begin{proof}

  We first prove the following lemma, greatly simplifying the remainder of the proof.

  \begin{lemma}
  \label{lem:c-to-n-map}
    Define $(s_{k}) = (2^{j})_{j=1}^{k}$, the sequence of the first $k$ integer powers of $2$. 
    Also, let $S_{c}$ be the set of such sequences all having exactly $c \in \mathbb{N}$ terms whose leading digit is $1$.
    If $(s_{m}) \in S_{c}$, then $c \log_{2}{10} < m < (c + 1) \log_{2}{10}$.
  \end{lemma}

  % TODO: Finish the argument 
  \begin{lemmaproof}

    Suppose that $(s_{m}) \in S_{c}$ for some $c \in \mathbb{N}$. We will show that $c \log_{2}{10} < m < (c + 1) \log_{2}{10}$.

    Observe that for $2^{q} \in \mathbb{R}$ to have a leading $1$, there must exist some $t \in \mathbb{N}$ such that 
    \[
      10^{t} \leq 2^{q} < 2 \cdot 10^{t} \implies t \log_{2}{10} \leq q < t \log_{2}{10} + 1.
    \]
    For fixed $t$, the interval $[t \log_{2}{10}, t \log_{2}{10} + 1)$ has length $1$, so $q$ can only take on, at most, one integer value. 

    Let $j_{1} < j_{2} < \cdots < j_{c} \leq m$ be the exponents of the terms of $(s_{m})$ that have a leading $1$. Therefore,
    each $j_{i} \in \mathbb{N}$ is contained in mutually exclusive unit intervals $[t_{i} \log_{2}{10}, t_{i} \log_{2}{10} + 1)$ and 
    $t_{1} < t_{2} < \cdots < t_{c}$. 
    Also, $t_{1} \geq 1$ since $2^{j_{1}} = 2^{4} = 16$ is the first integer power of $2$ with $1$ as its most significant digit.
    Because $(t_{i})_{i=1}^{c}$ is strictly increasing, we have $t_{i} \geq i$ for all $1 \leq i \leq c$.
    Thus, for all $1 \leq i \leq c$, yielding
    \[
      j_{i} \geq t_{i}\log_{2}{10} \geq i\log_{2}{10}.
    \]
    In particular, since $\log_{2}{10}$ is irrational, we have shown that 
    \[
      m \geq j_{c} > c \log_{2}{10} \implies m > c \log_{2}{10}.
    \]

    To establish the upperbound (i.e., $m < (c + 1) \log_{2}{10}$), we proceed in a similar fashion.
    Consider the next exponent $j_{c+1}$ wherein $2^{j_{c+1}}$ starts with a $1$. Clearly, $m < j_{c+1}$,
    and much of same reasoning above holds:
    \begin{enumerate}
      \item $\exists t_{c+1} \in \mathbb{N}.\, t_{c+1} \log_{2}{10} \leq j_{c+1} \leq t_{c+1} \log_{2}{10} + 1$
      \item $t_{1} < t_{2} < \cdots < t_{c} < t_{c+1}$
      \item $t_{c+1} \geq c + 1$.
    \end{enumerate}
    Now suppose to the contrary that $t_{c+1} \geq c + 2$. This would imply $j_{c+1} \not\in [t_{c+1} \log_{2}{10}, t_{c+1} \log_{2}{10} + 1)$,
    or equivalently, that this interval is devoid of an integer power of $2$. But, this gives us our desired contradiction:
    there exists $q \in \mathbb{N}$ such that $2^{q} \in [t_{c+1} \log_{2}{10}, t_{c+1} \log_{2}{10} + 1)$ contradicting 
    the fact that $j_{c+1}$ is the next exponent where $2^{j_{c+1}}$ starts with a $1$. Hence, we have
    $t_{c+1} = c + 1$, producing
    \[
      j_{c+1} < (c+1) \log_{2}{10} + 1.
    \]
    Finally, since $\log_{2}{10}$ is irrational, 

    \[
      m < j_{c + 1} < (c + 1) \log_{2}{10} \implies m < (c + 1) \log_{2}{10}.
    \]


  \end{lemmaproof}


    Turning our attention to the original problem, it is sufficient to show that 
    \[
    a_{n} \log_{10}{2} < n < (a_{n} + 1) \log_{10}{2}
    \]
    for all $n \in \mathbb{N}$ because
    \begin{align*}
      \log{2} - \frac{1}{n} < \frac{a_{n}}{n} < \log{2} &\iff n \log{2} - 1  < a_{n} < n \log{2} \\
                                                        &\iff \frac{n}{\log_{2}{10}} - 1 < a_{n} < \frac{n}{\log_{2}{10}} \\
                                                        &\iff a_{n} \log_{2}{10} < n < (a_{n} + 1) \log_{2}{10}. \\
    \end{align*}


    Let $n \in \mathbb{N}$ be arbitrary. Then, $(s_{n})$ has $a_{n}$ terms whose leading digit is 1 and $(s_{n}) \in S_{a_{n}}$.
    By Lemma $\ref{lem:c-to-n-map}$,
    \[
      a_{n} \log_{2}{10} < n < (a_{n} + 1)\log_{2}{10},
    \]
    as desired. Since $n$ is arbitrary, we can write 
    \[
      \forall n \in \mathbb{N}.\, \log{2} - \frac{1}{n} < \frac{a_{n}}{n} < \log{2}.
    \]
    By the \href{https://brilliant.org/wiki/squeeze-theorem/#statement-of-the-theorem-and-its-intuition}{Squeeze Theorem}, we immediately obtain
    \[
      \lim_{n \to \infty} \frac{a_{n}}{n} = \log{2} \approx 0.30103. 
    \]

  \end{proof}

\end{document}
